\documentclass[11pt]{article}
\usepackage[letterpaper,margin=0.8in]{geometry}
\usepackage[T1]{fontenc}
\usepackage[utf8]{inputenc}
\usepackage{lmodern}
\usepackage{array}
\usepackage{enumitem}
\usepackage[hidelinks]{hyperref}

\setlist[itemize]{leftmargin=1.2em,itemsep=2pt,topsep=2pt}
\setlength{\parskip}{0.5em}
\setlength{\parindent}{0pt}

\begin{document}

\begin{center}
{\Large \textbf{Mira Fellowship Application -- Written and Visual Materials}}\\
Piter Z. Garcia Bautista\\
Rochester, NY \quad $|$ \quad \href{mailto:pzg8794@rit.edu}{pzg8794@rit.edu} \quad $|$ \quad +1 (646) 288-3300
\end{center}

\section*{1. Idea Description, Impact Goals \& Timeline}

My project is \textbf{EQUITAS}, a trauma-informed and culturally responsive framework for building fairer diagnostic systems. It sits within my broader \textbf{Puzzle Plan} research direction: connecting bioinformatics, machine learning, education, and advocacy so that technical systems do not keep reproducing the same harms they claim to solve.

The core problem is that many diagnostic systems treat lived experience, trauma history, language difference, disability, and cultural context as noise instead of signal. When that happens, people who do not match the dominant clinical template are more likely to be dismissed, misread, or routed into the wrong interventions. This affects neurodivergent people, culturally and linguistically diverse communities, disabled patients, families navigating fragmented systems, and the professionals trying to support them with incomplete tools.

The deeper shift I want to create is a move from \textbf{extractive diagnosis} to \textbf{participatory diagnostic justice}. In that future, patient narrative is treated as data, not anecdote. Technical models are paired with accountability tools, community review, and decision pathways that are transparent enough to question and improve. The goal is not only better prediction. It is better recognition, better communication, and more humane action.

If EQUITAS succeeds, several things become possible that are still too rare today: earlier recognition of overlooked patterns, diagnostic tools that can be audited for exclusion, stronger bridges between health and education systems, and public-facing materials that help families and advocates understand what a system is doing and where it is failing.

During the Mira Fellowship, I would focus on a realistic six-month sprint:

\begin{itemize}
\item refine the problem framing and theory of change through structured stakeholder conversations;
\item turn the current research framework into a clear systems map, prototype toolkit, and public narrative;
\item identify one or two pilot pathways where EQUITAS could be tested with community-informed feedback;
\item build a partnership and fundraising roadmap for the next stage of work.
\end{itemize}

The Fellowship would help me move EQUITAS from a powerful interdisciplinary research direction into an actionable social venture with sharper language, stronger partnerships, and a clearer path to implementation.

\section*{2. Idea Impact Visual}

\begin{center}
\renewcommand{\arraystretch}{1.4}
\begin{tabular}{|>{\raggedright\arraybackslash}p{0.28\textwidth}|>{\centering\arraybackslash}p{0.07\textwidth}|>{\raggedright\arraybackslash}p{0.5\textwidth}|}
\hline
\textbf{Current state} & & \textbf{Future state with EQUITAS} \\
\hline
Symptoms are filtered through narrow norms; context is flattened or ignored. & $\rightarrow$ & Trauma history, language, disability, and lived experience are treated as meaningful parts of diagnostic interpretation. \\
\hline
Families and patients do not understand how a decision was made or how to challenge it. & $\rightarrow$ & Diagnostic pathways are legible, auditable, and paired with advocacy tools that support questions, appeals, and better next steps. \\
\hline
Health, education, and community systems respond separately to the same person. & $\rightarrow$ & Cross-sector tools create a shared picture of need, reducing misrouting and helping support plans line up across institutions. \\
\hline
AI and analytics risk scaling historic bias faster. & $\rightarrow$ & Fairness checks, community review, and reparative design are built in before scale, not added later as damage control. \\
\hline
\end{tabular}
\end{center}

This visual shows the world I am trying to help build: one where diagnostic systems become more accurate because they become more human. The change is not only technical. It is relational and institutional. Patients gain more dignity and clarity. Families gain tools for action. Clinicians and educators gain better context, not just more scores. Communities gain a way to question and reshape the systems that classify them. I want EQUITAS to make fairer diagnosis feel normal, expected, and operational rather than exceptional.

\section*{3. Resource Analysis}

\textbf{What I have:} I already have an unusual interdisciplinary foundation for this work. I bring training across computer science, data science, computational biology, and inclusive education; professional experience in healthcare AI and data systems; current research on reproducible machine learning evaluation; and an emerging body of writing and framework design around EQUITAS and the Puzzle Plan. I also bring lived experience with delayed and fragmented diagnosis, which gives this work urgency and keeps it grounded in consequences rather than abstraction.

\textbf{What I need:} The missing pieces are not only technical. I need founder-level coaching, stronger community and practitioner partnerships, design support for turning frameworks into usable tools, legal and policy guidance around responsible deployment, and protected time to move beyond academic framing into venture framing. I also need a small circle of clinical, advocacy, and systems-change advisors who can stress-test the work from different angles. Financially, micro-grants and early-stage support would materially accelerate stakeholder engagement, prototype development, and pilot preparation.

\textbf{Milestones:}

\begin{itemize}
\item \textbf{0--12 months:} finalize theory of change, complete stakeholder interviews, produce a prototype toolkit and systems map, and establish an initial advisory circle.
\item \textbf{12--24 months:} run one or two community-informed pilots, document lessons, publish open explanatory materials, and secure anchor partnerships or seed support.
\item \textbf{24--36 months:} formalize the operating structure, expand into a repeatable partnership model, and build an evaluation framework that measures both technical performance and dignity-centered outcomes.
\end{itemize}

Several assumptions still need testing. I need to learn which parts of the framework communities and institutions find most useful first, where resistance shows up, and how much can be operationalized without losing the core justice commitments. That is exactly why Mira is a fit. I do not need passive encouragement. I need rigorous, iterative pressure that helps the idea mature.

\section*{4. Why You?}

This work matters to me because I know what it feels like to be intelligible to a system only in fragments. My own diagnostic journey included years of delayed recognition, partial explanations, and moments when my experience was treated as secondary to what institutions expected someone like me to look or sound like. That taught me something I have not been able to unsee: many systems fail not because people are too complex, but because the systems are too narrow.

Professionally, I have spent the last several years moving across technical and human-centered domains that are usually kept apart. I have worked in healthcare AI and data settings, completed graduate work in computer science and data science, built research workflows in machine learning and bioinformatics, and simultaneously pursued graduate study in teaching, disability, and inclusive practice. That combination has shaped how I see the problem. I do not want to build a better model inside the same old frame. I want to redesign the frame itself.

What makes me particularly suited to this work is the combination of technical fluency, interdisciplinary range, and lived proximity to the stakes. I can speak with researchers about models, pipelines, evaluation, and reproducibility. I can also speak with educators, families, and advocates about trust, language, access, and why a technically elegant system can still do real harm. EQUITAS exists because I have lived at that intersection long enough to know that the gap between those worlds is itself part of the problem.

I am not applying because this is an interesting side project. I am applying because this is the work I want to spend my life building.

\section*{5. Why Now?}

This work must begin now because diagnostic systems are entering a dangerous acceleration phase. AI is moving from research environments into real decision pathways across healthcare, education, and public systems. The tools are getting more powerful at exactly the same moment institutions are under pressure to act faster, automate more, and justify decisions with data. If we do not intervene now, we risk hard-coding old inequities into the next generation of supposedly intelligent systems.

There is also an unusual window of opportunity. Multimodal health data is more available. Public conversation about trustworthy AI is more developed. Communities harmed by exclusionary systems are naming the problem with greater clarity and force. At the same time, many institutions still lack practical frameworks that connect fairness, lived experience, implementation, and accountability. EQUITAS is my attempt to help fill that gap.

For me personally, this is also the moment when several strands of my life have finally converged. My graduate research, healthcare technology experience, bioinformatics training, and work in inclusive education are no longer separate tracks. They are pointing toward the same question: how do we build systems that can recognize people more fully and act more justly? I now have enough clarity to name the problem, enough evidence to show it is real, and enough momentum to begin building something durable.

If this work is delayed, the opportunity cost is not only personal. The cost is that more systems will be scaled before they are questioned, and more people will continue to be misread by institutions that claim neutrality. I want to use this moment to move from research and articulation into structure, partnership, and implementation.

\end{document}